\section*{Capítulo \ref{ch:b2:himalaya-andes} --- Cinturões de colisão e cinturões andinos comparados}

\begin{solution}{exo:b2:himalaya-andes:1}
(a) $r = h\rho_c/(\rho_m - \rho_c) = 3.0\times2700/600 = \qty{13.5}{km}$. (b) $T = 35 + 3.0 +
13.5 = \qty{51.5}{km}$, ou $35 + 5.5\times3.0$. (c) $h = \Delta T(\rho_m - \rho_c)/\rho_m =
10\times600/3300 = \qty{1.8}{km}$.
\end{solution}

\begin{solution}{exo:b2:himalaya-andes:2}
(a) Antefossa indo-gangética; Cavalgamento Frontal Principal; Sub-Himalaia (Siwalik); Cavalgamento
Limítrofe Principal; Baixo Himalaia; Cavalgamento Central Principal; sequência do Alto Himalaia; Sistema de
\omterm{def:b2:alpine-orogen:decollement}{Descolamento} do Sul do Tibete; sequência do Himalaia Tetiano; depois a sutura Indo--Tsangpo e a crosta
asiática. (b) O Sistema de \omterm{def:b2:alpine-orogen:decollement}{Descolamento} do Sul do Tibete. (c) O grau aumenta para o norte: \omterm{def:b2:alpine-orogen:foreland}{molassa}
não metamorfizada, Baixo Himalaia de grau médio, sequência do Alto Himalaia de alto grau e parcialmente
fundida; acima do STDS, ele cai bruscamente para sedimentos tetianos não metamorfizados ou fracamente
metamorfizados, pois a falha normal removeu as rochas intermediárias.
\end{solution}

\begin{solution}{exo:b2:himalaya-andes:3}
(a) $196.5\ \text{m}/9200\ \text{yr} = \qty{0.02136}{m}$ por ano, cerca de \qty{21.4}{mm} por ano.
(b) É \qty{20}{\percent} a mais que os \qty{17.8}{mm} por ano do GPS, mas da mesma ordem:
uma taxa é a média sobre nove mil anos, a outra sobre vinte.
(c) Quase todo o encurtamento através da cadeia é absorvido na frente: o MBT, o MCT e
as outras falhas ao norte do MFT estão, nesta longitude, praticamente inativos; todos se enraízam
no mesmo MHT, que aflora no MFT.
\end{solution}

\begin{solution}{exo:b2:himalaya-andes:4}
(a) Colisão: a sutura marca o oceano desaparecido. (b) Tipo andino: um arco precisa de uma laje
que desce. (c) Colisão, tipicamente: fundidos crustais de uma crosta espessada e aquecida (eles também
ocorrem em crosta andina espessada, mas ali dominam os \omterm{def:b2:magma-emplacement:magma}{magmas} de arco). (d) Tipo andino.
(e) Os dois: as Siwalik e o antepaís subandino. (f) Os dois: o MHT no Himalaia, a
interface de subducção nos Andes.
\end{solution}

\begin{solution}{exo:b2:himalaya-andes:5}
(a) $17.8/45 = \qty{40}{\percent}$ no leste, $20.5/45 = \qty{46}{\percent}$ no oeste:
cerca da metade. (b) Para o espessamento e as falhas do Tibete e do resto da Ásia, e para
a extrusão de blocos para leste ao longo das grandes falhas transcorrentes. (c) $45 - 17.8 =
\qty{27.2}{mm}$ por ano, ou \qty{272}{km} em \qty{10}{Myr}; $45 - 20.5 = \qty{24.5}{mm}$ por
ano, ou \qty{245}{km}.
\end{solution}

\begin{solution}{exo:b2:himalaya-andes:6}
(a) $r = 4.5\times3.7 = \qty{16.65}{km}$; $\Delta F = 2700\times9.8\times3700\times(35\,000 +
1850 + 8325) = \qty{4.42e12}{N.m^{-1}}$. (b) Tibete: \qty{6.12e12}{N.m^{-1}}; o Altiplano
tem \qty{72}{\percent} desse valor. (c) Com $r = 4.5h$, $\Delta F = \rho_c g\,(Hh + 2.75h^2)$: para
$h \ll H$, o termo $Hh$ domina e $\Delta F$ é proporcional a $h$; para um planalto alto,
o termo em $h^2$ cresce mais depressa. Cada quilômetro a mais de altura custa cada vez mais empurrão nas
bordas, até que as forças disponíveis já não conseguem elevar o planalto; ele então se espalha
em vez de subir.
\end{solution}

\begin{solution}{exo:b2:himalaya-andes:7}
(a) $2700\times9.8\times8848.86/2 = \qty{1.17e8}{Pa} = \qty{117}{MPa}$. (b) A
\qty{100}{MPa}: $0.85\times100 = \qty{85}{MPa}$; a \qty{200}{MPa}: $0.85\times200 = 50 +
0.6\times200 = \qty{170}{MPa}$, os dois ramos concordam. (c) $h = 2\sigma_d/(\rho g) =
2\times1.70\times10^8/(2700\times9.8) = \qty{12.8}{km}$: uma ordem de grandeza, da ordem
dos cumes mais altos, não muito acima deles.
\end{solution}

\begin{solution}{exo:b2:himalaya-andes:8}
(a) $11\times111.19\times\cos\ang{30} = \qty{1059}{km}$. (b) $9\times10^{-3}/1.059\times10^{6}
= \num{8.5e-9}$ por ano (entre \num{6.6e-9} e \num{10.4e-9} para 7 e \qty{11}{mm} por ano),
isto é, $8.5\times10^{-9}/3.156\times10^{7} = \qty{2.7e-16}{s^{-1}}$. (c) Os dois intervalos
se sobrepõem: $(12.4 \pm 4.7) \times 10^{-9}$ vai de 7.7 a \num{17.1e-9}. O sul do planalto, ao norte do
Nepal, se estira tão depressa quanto o cinturão como um todo, ou um pouco mais.
\end{solution}

\begin{solution}{exo:b2:himalaya-andes:9}
(a) $0.9/6.6 = \qty{14}{\percent}$. (b) $M_0 = 10^{1.5\times8.5 + 9.1} = 10^{21.85} =
\qty{7.1e21}{N.m}$. (c) $7.1\times10^{21}/6.6\times10^{19} = 107$: um terremoto assim a cada
110 anos, aproximadamente, sob o Nepal, muito mais do que a história registra; ou terremotos muito maiores
estão por vir, ou parte do déficit é liberada de outra forma.
\end{solution}

\begin{solution}{exo:b2:himalaya-andes:10}
(a) Tome-se a profundidade de compensação na base da crosta do Altiplano, $61.5 - 3.7 =
\qty{57.8}{km}$ abaixo do nível do mar. A crosta da Puna chega a $50.5 - 4.4 = \qty{46.1}{km}$, de modo que a
coluna da Puna contém \qty{11.7}{km} de manto acima dessa profundidade: $2800\times61.5 =
\rho_P\times50.5 + 3300\times11.7$. (b) $\rho_P = (172\,200 - 38\,610)/50.5 =
\qty{2645}{kg.m^{-3}}$. (c) A diferença, \qty{155}{kg.m^{-3}}, fica no extremo inferior dos
150 a \qty{200}{kg.m^{-3}} citados: só a crosta exigiria uma crosta da Puna excepcionalmente leve.
Uma coluna mais leve abaixo dela o faz com mais facilidade: um manto litosférico removido por delaminação,
substituído por astenosfera quente e leve.
\end{solution}

\begin{solution}{exo:b2:himalaya-andes:11}
(a) $g = 105/\tan\ang{30} = \qty{181.9}{km}$. (b) $g = 90/\tan\ang{30} + 15/\tan\ang{4} =
155.9 + 214.5 = \qty{370.4}{km}$: o arco se desloca \qty{188.5}{km} para o interior. (c) A laje nunca
atinge \qty{105}{km} sob a placa ao longo desses \qty{300}{km}, e nenhuma cunha mantélica quente fica
acima dela: nenhum \omterm{def:b2:magma-emplacement:magma}{magma} é produzido e o arco morre até que a laje volte a afundar, mais a leste.
\end{solution}

\begin{solution}{exo:b2:himalaya-andes:12}
(a) \qty{21}{mm} por ano são \qty{21}{km} por milhão de anos: $21\times37.5 =
\qty{787.5}{km^3}$ por quilômetro de cadeia por milhão de anos. (b) $300/787.5 =
\qty{38}{\percent}$. (c) No espessamento da cadeia e do sul do Tibete, na crosta
levada para baixo com a Índia sob o Tibete, e no material que o MHT transfere para o norte.
\end{solution}

\begin{solution}{pb:b2:himalaya-andes:1}
\textbf{1.} Sete amostras de \ang{28.0} a \ang{37.0}~N: $34\,550/7 = \qty{4936}{m}$, $h_T
= \qty{4.9}{km}$.
\textbf{2.} Cinco amostras de \ang{-69} a \ang{-65}: $18\,650/5 = \qty{3730}{m}$, $h_A =
\qty{3.7}{km}$, o valor citado para o Altiplano; a Puna, mais ao sul, é mais alta.
\textbf{3.} Tibete: pelo menos $9\times111.19 = \qty{1001}{km}$; o Altiplano: pelo menos $4\times
104.5 = \qty{418}{km}$.
\textbf{4.} $(5000 - 70)/(1.5\times111\,190) = 0.0296$: \qty{3.0}{\percent}.
\textbf{5.} O perfil andino começa no fundo oceânico e na fossa, quilômetros abaixo do
nível do mar, a oeste do planalto; o tibetano começa numa planície continental perto do nível do mar,
e termina na alta bacia do Tarim.
\textbf{6.} $r = 4.5\times4.9 = \qty{22.05}{km}$; $T_T = 35 + 4.9 + 22.05 = \qty{62.0}{km}$.
\textbf{7.} $r = 4.5\times3.7 = \qty{16.65}{km}$; $T_A = 35 + 3.7 + 16.65 = \qty{55.4}{km}$.
\textbf{8.} Tibete: \qty{62.0}{km} contra 70 a \qty{75}{km}, \qty{10.5}{km} abaixo do valor
médio \qty{72.5}{km} (\qty{14}{\percent}); o Altiplano: \qty{55.4}{km} contra 59 a
\qty{64}{km}, \qty{6.1}{km} abaixo de \qty{61.5}{km} (\qty{10}{\percent}). O Tibete está mais longe.
\textbf{9.} $72.5 - 35 = 4.9\times3300/(3300 - \rho_c)$: $3300 - \rho_c = 16\,170/37.5 =
431.2$, $\rho_c = \qty{2869}{kg.m^{-3}}$.
\textbf{10.} $26.5 = 3.7\times3300/(3300 - \rho_c)$: $3300 - \rho_c = 460.8$, $\rho_c =
\qty{2839}{kg.m^{-3}}$.
\textbf{11.} Parte da carga é sustentada pela flexão da placa, e não por uma raiz local; a
crosta normal pode não ter tido \qty{35}{km} de espessura; o manto abaixo pode ser mais leve ou mais denso
do que o suposto; o Moho sísmico pode não ser a base da crosta no sentido da densidade.
\textbf{12.} A base da crosta do Tibete fica a $72.5 - 4.9 = \qty{67.6}{km}$, a da Índia a
\qty{37.5}{km}: $t = \qty{30.1}{km}$. $\Delta g = 2\pi\times6.674\times10^{-11}\times370\times
30\,100 = \qty{4.67e-3}{m.s^{-2}}$, isto é, \qty{-467}{mGal} em relação à Índia, perto da
diferença observada de \qty{460}{mGal}.
\textbf{13.} $W_0 = 440\times60/35 = \qty{754.3}{km}$; $S = \qty{314.3}{km}$.
\textbf{14.} $(314.3 - 250)/314.3 = \qty{20}{\percent}$, o déficit encontrado nos Andes
Centrais.
\textbf{15.} $0.55\times314.3 = \qty{173}{km}$, menos que os \qty{250}{km} das \omterm{def:b2:alpine-orogen:balanced}{seções
balanceadas}, que seriam então mais do que suficientes.
\textbf{16.} A cobertura encurtada acima do \omterm{def:b2:alpine-orogen:decollement}{descolamento} deve ser compensada por um encurtamento do
embasamento abaixo dela, mais a oeste, pois os dois devem perder o mesmo comprimento através do cinturão; o encurtamento
do embasamento espessa a crosta sob a Cordilheira Oriental e o planalto.
\textbf{17.} $314.3/30 = \qty{10.5}{mm}$ por ano, $10.5/72 = \qty{15}{\percent}$ da
convergência.
\textbf{18.} Ela é consumida pela subducção: a placa de Nazca afunda sob o continente.
\textbf{19.} $r = \qty{22.05}{km}$; $\Delta F = 2700\times9.8\times4900\times(35\,000 + 2450
+ 11\,025) = \qty{6.28e12}{N.m^{-1}}$.
\textbf{20.} $2700\times9.8\times3700\times45\,175 = \qty{4.42e12}{N.m^{-1}}$; a razão é de
0.70.
\textbf{21.} $6.28\times10^{12}/10^{5} = \qty{6.3e7}{Pa}$, cerca de \qty{63}{MPa}.
\textbf{22.} $9\times10^{-3}/1.06\times10^{6} = \num{8.5e-9}$ por ano.
\textbf{23.} O empurrão da Índia atua na direção norte--sul e é equilibrado, no alto planalto, pelo
próprio peso do planalto: na direção norte--sul o planalto já não pode se espessar, de modo que a tensão vertical,
o peso, é a maior; a menor tensão é leste--oeste, em direção às bordas livres da Ásia,
e o planalto se estende nessa direção, em riftes norte--sul.
\textbf{24.} \textbf{\qty[text-series-to-math=true]{62}{km}}, cerca de \qty{86}{\percent} dos
\qty{72.5}{km} medidos.
\end{solution}
